% !TEX program = pdflatex
% ------------------------------------------------------------
% MATH 431 Homework Template (plug-and-chug friendly)
%
% Goals:
% - Same clean header style as your MATH 521 template.
% - Easy “Problem -> Solution” workflow.
% - Lightweight (no tcolorbox); uses mdframed for optional boxes.
% - Good for probability homework: sample space, probability measure, events.
%
% Quick use:
%   1) Edit the Metadata block.
%   2) For each exercise, use \problem{<section>.<number>} and paste the prompt.
%   3) Write your solution in the solution environment.
% ------------------------------------------------------------

\documentclass[12pt]{article}

% =========================
% 0) Metadata (EDIT ME)
% =========================
\newcommand{\CourseCode}{MATH 431}
\newcommand{\CourseTitle}{Introduction to Probability}
\newcommand{\CourseSection}{LEC 002} % change to 001 or 003 if needed
\newcommand{\Semester}{Spring 2026}
\newcommand{\Instructor}{Sarah Tammen} % LEC 001: Ander Aguirre | LEC 002: Sarah Tammen | LEC 003: Nicholas Derr
\newcommand{\MeetingInfo}{MWF 9:55--10:45, 6102 Sewell Social Sciences} % LEC 001 default; update if in 002/003
\newcommand{\HomeworkNumber}{5}
\newcommand{\YourName}{Alejandro De La Torre}
\newcommand{\DueDate}{\today} % or e.g. February 2, 2026

% =========================
% 1) Core math tools
% =========================
\usepackage{amsmath, amssymb, amsthm}

% =========================
% 2) Lists and spacing
% =========================
\usepackage{enumitem}
\setlist[enumerate,1]{label=\textbf{(\alph*)}, leftmargin=2em, itemsep=0.25em}
\setlist[itemize]{leftmargin=1.5em, itemsep=0.25em}
\setlength{\parskip}{0.35em}
\setlength{\parindent}{0pt}

% =========================
% 3) Page geometry + header/footer
% =========================
\usepackage{geometry}
\geometry{margin=1in}

\usepackage{fancyhdr}
\pagestyle{fancy}
\fancyhf{}
\lhead{\CourseCode\ --- \CourseSection, \Semester}
\rhead{Homework \HomeworkNumber}
\cfoot{\thepage}
\renewcommand{\headrulewidth}{0.4pt}
\setlength{\headheight}{15pt}
\addtolength{\topmargin}{-2pt}

% =========================
% 4) Links (subtle)
% =========================
\usepackage[hidelinks]{hyperref}
\setcounter{tocdepth}{1}

% =========================
% 5) Figures (optional)
% =========================
\usepackage{graphicx}
\graphicspath{{images/}}

% =========================
% 6) Lightweight boxes (optional)
% =========================
\usepackage{mdframed}
\newmdenv[
  skipabove=0.6\baselineskip,
  skipbelow=0.6\baselineskip,
  linewidth=0.6pt,
  roundcorner=4pt,
  innerleftmargin=10pt,
  innerrightmargin=10pt,
  innertopmargin=8pt,
  innerbottommargin=8pt
]{boxnote}

% =========================
% 7) Probability / notation macros
% =========================
\newcommand{\R}{\mathbb{R}}
\newcommand{\N}{\mathbb{N}}
\newcommand{\Z}{\mathbb{Z}}

\newcommand{\OmegaSpace}{\Omega}
\newcommand{\FField}{\mathcal{F}}
% Probability measure in case it is relevant or want to distinguish it from P(A)
%\newcommand{\PMeasure}{\mathbb{P}} 
\newcommand{\E}{\mathbb{E}}

\newcommand{\given}{\,\middle|\,}
\newcommand{\ind}{\mathbf{1}}

% Common “defined as” symbol
\newcommand{\defeq}{\vcentcolon=}

% =========================
% 8) Problem command (TOC + PDF bookmarks)
% Usage:
%   \problem{<label>}
%   \problem[Short title]{<label>}
% Example:
%   \problem{1.4}
%   \problem[Sample space]{1.4}
% =========================
\makeatletter
\newcommand{\problem}[2][]{%
  \def\prob@title{Problem #2\if\relax\detokenize{#1}\relax\else\ (\textit{#1})\fi}%
  \phantomsection%
  \pdfbookmark[1]{\prob@title}{prob#2}%
  \addcontentsline{toc}{section}{\prob@title}%
  \section*{\prob@title}%
  \vspace{-0.25em}\hrule\vspace{0.8em}%
}
\makeatother

% =========================
% 9) Solution environment
% =========================
\newenvironment{solution}{%
  \noindent\textbf{Solution.}\ \ignorespaces
}{\par}

% Optional scratch / planning block
\newenvironment{scratch}{%
  \begin{boxnote}\textbf{Scratch / Setup.}\ \small
}{\end{boxnote}}

% Final answer command: expects math only (no $$ or \[ \])
\newcommand{\finalanswer}[1]{%
  \par\medskip
  \noindent\textbf{Final:}\ \fbox{$#1$}\par
} % AGAIN: MATH ONLY INSIDE THE BOX; if word answer, use \text{} inside or use \fbox{TEXT} or \framebox

% =========================
% 10) Title block
% =========================
\title{Homework \HomeworkNumber}
\author{\YourName \\
\CourseCode\ --- \CourseTitle \\
\CourseSection \\
Instructor: \Instructor}
\date{Due: \DueDate}

\begin{document}
\maketitle

% \section*{Collaboration / Resources}
% (Optional) Write “I worked alone” or list collaborators/resources.

\tableofcontents
\bigskip
\hrule
\bigskip
\newpage

% ============================================================
% Problems (duplicate blocks as needed)
% ============================================================


\problem{3.5}
Suppose that the discrete random variable $X$ has cumulative distribution function given by
\[
F(x)=
\begin{cases}
0, & \text{if } x < 1,\\
\frac{1}{3}, & \text{if } 1 \le x < \frac{4}{3},\\
\frac{1}{2}, & \text{if } \frac{4}{3} \le x < \frac{3}{2},\\
\frac{3}{4}, & \text{if } \frac{3}{2} \le x < \frac{9}{5},\\
1, & \text{if } x \ge \frac{9}{5}.
\end{cases}
\]
Find the possible values and the probability mass function $X$.

\begin{solution}
\begin{scratch}
The c.d.f. jumps at the support points. Each point mass equals the jump size:
$P(X=x)=F(x)-\lim_{t\uparrow x}F(t)$.
\end{scratch}

The jump points are $1,\frac{4}{3},\frac{3}{2},\frac{9}{5}$, so these are the
possible values. The masses are
\[
P(X=1)=\frac{1}{3}-0=\frac{1}{3},\qquad
P\!\left(X=\frac{4}{3}\right)=\frac{1}{2}-\frac{1}{3}=\frac{1}{6},
\]
\[
P\!\left(X=\frac{3}{2}\right)=\frac{3}{4}-\frac{1}{2}=\frac{1}{4},\qquad
P\!\left(X=\frac{9}{5}\right)=1-\frac{3}{4}=\frac{1}{4}.
\]
These sum to $1$.

\finalanswer{X\in\left\{1,\frac{4}{3},\frac{3}{2},\frac{9}{5}\right\},\;
P(X=1)=\frac{1}{3},\;
P\!\left(X=\frac{4}{3}\right)=\frac{1}{6},\;
P\!\left(X=\frac{3}{2}\right)=\frac{1}{4},\;
P\!\left(X=\frac{9}{5}\right)=\frac{1}{4}}

\end{solution}

\newpage

\problem{3.7}
Suppose that the continuous random variable $X$ has cumulative distribution function given by
\[
F(x)=
\begin{cases}
0, & \text{if } x < \sqrt{2},\\
x^2 - 2, & \text{if } \sqrt{2} \le x < \sqrt{3},\\
1, & \text{if } \sqrt{3} \le x.
\end{cases}
\]
\begin{enumerate}[label=(\alph*)]
\item Find the smallest interval $[a,b]$ such that of $P(a \le X \le b) = 1$.
\item Find $P(X = 1.6)$.
\item Find $P(1 \le X \le \frac{3}{2})$.
\item Find the probability density function of $X$.
\end{enumerate}

\begin{solution}
\begin{scratch}
Support is where $F$ rises from 0 to 1, i.e., $[\sqrt{2},\sqrt{3}]$.
Since $X$ is continuous, point probabilities are 0. The pdf is $f=F'$ on the
interval where $F$ is differentiable.
\end{scratch}

\begin{enumerate}[label=(\alph*)]
\item The smallest interval with probability 1 is
\[
[a,b]=\bigl[\sqrt{2},\sqrt{3}\bigr].
\]
\item Since $X$ is continuous,
\[
P(X=1.6)=0.
\]
\item
\[
P\!\left(1\le X\le \frac{3}{2}\right)=F\!\left(\frac{3}{2}\right)-F(1)
\]
and $1<\sqrt{2}$, so $F(1)=0$ and
\[
F\!\left(\frac{3}{2}\right)=\left(\frac{3}{2}\right)^2-2=\frac{1}{4}.
\]
Thus $P(1\le X\le 3/2)=1/4$.
\item
\[
f_X(x)=F'(x)=
\begin{cases}
2x, & \sqrt{2}<x<\sqrt{3},\\
0, & \text{otherwise.}
\end{cases}
\]
\end{enumerate}

\finalanswer{\text{(a)}\ [\sqrt{2},\sqrt{3}],\ \text{(b)}\ 0,\ \text{(c)}\ \frac{1}{4},\ 
\text{(d)}\ f_X(x)=\begin{cases}2x,&\sqrt{2}<x<\sqrt{3}\\0,&\text{otherwise}\end{cases}}

\end{solution}

\newpage

\problem{3.9}
Let $X$ be the random variable from Exercise 3.3.
\begin{enumerate}[label=(\alph*)]
\item Find the mean of $X$.
\item Compute $E[e^{2X}]$.
\end{enumerate}

\begin{solution}
\begin{scratch}
From Exercise 3.3, $f(x)=3e^{-3x}$ for $x>0$ (and 0 otherwise), so
$X\sim\text{Exp}(3)$. Use
$E[X]=\int_0^\infty x f(x)\,dx$ and
$E[e^{2X}]=\int_0^\infty e^{2x}f(x)\,dx$.
\end{scratch}

\begin{enumerate}[label=(\alph*)]
\item
\[
E[X]=\int_{0}^{\infty} x\cdot 3e^{-3x}\,dx=\frac{1}{3}.
\]
\item
\[
E[e^{2X}]=\int_{0}^{\infty} e^{2x}\cdot 3e^{-3x}\,dx
=3\int_{0}^{\infty}e^{-x}\,dx=3.
\]
\end{enumerate}

\finalanswer{\text{(a)}\ E[X]=\frac{1}{3},\quad \text{(b)}\ E[e^{2X}]=3}

\end{solution}

\newpage

\problem{3.10}
Let $X$ have probability mass function
\[
P(X = -1) = \frac{1}{2}, \quad P(X = 0) = \frac{1}{3}, \quad \text{and } P(X = 1) = \frac{1}{6}.
\]
Calculate $E[|X|]$ using the approaches in (a) and (b) below.
\begin{enumerate}[label=(\alph*)]
\item First find the probability mass function of the random variable $Y = |X|$ and using that compute $E[|X|]$.
\item Apply formula (3.24) with $g(x) = |x|$.
\end{enumerate}

\begin{solution}
\begin{scratch}
Let $Y=|X|$. Then $Y\in\{0,1\}$ with
$P(Y=0)=P(X=0)=1/3$ and
$P(Y=1)=P(X=1)+P(X=-1)=2/3$.
\end{scratch}

\begin{enumerate}[label=(\alph*)]
\item The p.m.f. of $Y$ is
\[
P(Y=0)=\frac{1}{3},\qquad P(Y=1)=\frac{2}{3}.
\]
Thus
\[
E[|X|]=E[Y]=0\cdot\frac{1}{3}+1\cdot\frac{2}{3}=\frac{2}{3}.
\]
\item Using (3.24) with $g(x)=|x|$,
\[
E[|X|]=\sum_x |x|P(X=x)
 =1\cdot\frac{1}{2}+0\cdot\frac{1}{3}+1\cdot\frac{1}{6}=\frac{2}{3}.
\]
\end{enumerate}

\finalanswer{E[|X|]=\frac{2}{3}}

\end{solution}

\newpage

\problem{3.47}
Choose a point uniformly at random from the triangle with vertices $(0, 0)$, $(30, 0)$, and $(30, 20)$. Let $(X, Y)$ be the coordinates of the chosen point.
\begin{enumerate}[label=(\alph*)]
\item Find the cumulative distribution function of $X$.
\item Use part (a) to find the density function of $X$.
\item Use part (b) to find the expectation of $X$.
\end{enumerate}

\begin{solution}
\begin{scratch}
The triangle has area $\frac{1}{2}\cdot 30\cdot 20=300$.
The upper boundary line from $(0,0)$ to $(30,20)$ is $y=\frac{2}{3}x$.
For $0\le x\le 30$, the area with $X\le x$ is
$\int_0^x \frac{2}{3}t\,dt=\frac{1}{3}x^2$.
\end{scratch}

\begin{enumerate}[label=(\alph*)]
\item
\[
F_X(x)=
\begin{cases}
0, & x<0,\\[4pt]
\dfrac{x^2}{900}, & 0\le x<30,\\[6pt]
1, & x\ge 30.
\end{cases}
\]
\item Differentiate:
\[
f_X(x)=
\begin{cases}
\dfrac{x}{450}, & 0\le x<30,\\[6pt]
0, & \text{otherwise.}
\end{cases}
\]
\item
\[
E[X]=\int_{0}^{30} x\,f_X(x)\,dx
=\int_{0}^{30} x\cdot \frac{x}{450}\,dx
=\frac{1}{450}\cdot \frac{30^3}{3}=20.
\]
\end{enumerate}

\finalanswer{\text{(a)}\ F_X(x)=\begin{cases}0,&x<0\\x^2/900,&0\le x<30\\1,&x\ge30\end{cases}}
\finalanswer{\text{(b)}\ f_X(x)=\begin{cases}x/450,&0\le x<30\\0,&\text{otherwise}\end{cases}}
\finalanswer{\text{(c)}\ E[X]=20}

\end{solution}

\newpage

\problem{3.48}
Pick a random point uniformly inside the triangle with vertices $(0, 0)$, $(2, 0)$ and $(0, 1)$. Compute the expectation of the distance of this point to the $y$-axis.

\begin{solution}
\begin{scratch}
The distance to the $y$-axis equals the $x$-coordinate. The triangle has area
$\frac{1}{2}\cdot 2\cdot 1=1$ and upper boundary $y=1-\frac{x}{2}$ for
$0\le x\le 2$.
\end{scratch}

Let $R$ be the distance to the $y$-axis; then $R=X$. Thus
\[
E[R]=\frac{1}{\text{Area}}\int_{0}^{2} x\left(1-\frac{x}{2}\right)\,dx
=\int_{0}^{2}\left(x-\frac{x^2}{2}\right)\,dx
=\left[\frac{x^2}{2}-\frac{x^3}{6}\right]_{0}^{2}
=\frac{2}{3}.
\]

\finalanswer{\frac{2}{3}}

\end{solution}


\end{document}
